% MT012 / explanation — canonical editable content.
先给每个符号一个位置：第\(i\)条训练样本为\(x_i\)，标签\(y_i\in\{-1,+1\}\)；\(w\)是特征权重，\(b\)是偏置；\(f(x_i)=w^Tx_i+b\)是分类分数。\(y_if(x_i)>0\)表示分数方向和真实类别一致。


原来L2对权重平方求和，L1则对绝对值求和：\(\|w\|_1=\sum_j|w_j|\)。只更换正则项，Logistic数据损失不变，所以得到英文答案中的目标。\(\arg\min\)表示寻找使整项最小的\(w,b\)；这里沿用不惩罚\(b\)的约定，\(C>0\)控制权重惩罚。


L1的特点是容易让某些权重恰好为0。一个权重为0时，对应特征无论怎样改变，都不再贡献分数，这就是稀疏（sparse）模型。比如\(w=(0,2)\)，分数\(2x_2+b\)只看纵坐标。


回到题图，红叉与蓝圈主要沿纵向分开，横向位置对区分类别帮助较小。因此合理的稀疏草图是\textbf{水平分隔线}：取\(w_1=0\)，边界满足



\[
w_2x_2+b=0\quad\Longrightarrow\quad x_2=-b/w_2\quad(w_2\ne0).
\]


读图时在上下两群之间画水平线，并说明“保留\(x_2\)，舍去\(x_1\)”。原图没有足够数值来确定唯一高度，不需要编一个精确截距。L1是鼓励稀疏，不是保证所有数据都得到轴平行边界；弱正则时两个系数都可能保留。


选读：平方惩罚在权重接近0时导数也接近0；绝对值在0两侧仍有固定大小的斜率，在0处形成尖点。这使一段范围内的数据拟合收益不足以抵消L1代价，最优权重就停在0，详见\StudyRef{Answers}{mt019}{MT019}。


\textbf{出处：} 2020B期中Quiz Q12，原卷第6页（10分）


\textbf{答案依据：} 2020B期中 Q12 随卷／配套参考答案；条件与解释经复核。
